% --------------------------------------------------------------------
%  SMALL COMMANDS
% --------------------------------------------------------------------
%  \term for a defined term, \cmd for inline code, \chapref and
%  \secref for cross-references, and the utpsteps numbered-step list.
%
%      \term{overfitting}          a term being defined
%      \cmd{git status}            inline code, in any language
%      \chapref{ch:trees}          "Chapter 12"
%      \secref{sec:pruning}        "../labs/layout.tex.3"
%      \chapterguide{prior}{outcomes}   the box that opens a chapter
%      \begin{utpsteps} ... \end{utpsteps}  a numbered Step 1/Step 2 list
%
%  The second group is the shorthand the chapters use in running text:
%
%      \figref{fig:tree}           "Figure 4.2"
%      \chapsref{ch:a}{ch:b}       "Chapters 3--5"
%      \lib{pandas}                a library name
%      \cxg{x}, \cxo{x}, \cxr{x}   maths coloured green/orange/red
%      \given                      the bar in P(a \given b)
%      \coretag, \latertag         read-now / come-back-to markers
%      \kagglesource{url}{name}    the dataset a lab is built on

\newcommand{\term}[1]{\textbf{\color{UTPNavy}#1}}
\newcommand{\cmd}[1]{\texttt{#1}}
\newcommand{\chapref}[1]{Chapter~\ref{#1}}
\newcommand{\secref}[1]{Section~\ref{#1}}
\newcommand{\chapterguide}[2]{%
  \begin{tcolorbox}[utpplain,breakable,colback=white,colframe=UTPNavy,boxrule=0.8pt,
    arc=2mm,left=2.5mm,right=2.5mm,top=2mm,bottom=2mm,
    title=\utptitleglance,fonttitle=\bfseries\color{white},colbacktitle=UTPNavy]
    \textbf{Before you start:} #1\par
    \textbf{By the end, you can:} #2
  \end{tcolorbox}}
\newenvironment{utpsteps}
  {\begin{enumerate}[label=\textcolor{UTPBlue}{\bfseries Step~\arabic*:},leftmargin=4.5em,itemsep=0.38em]}
  {\end{enumerate}}

\newcommand{\figref}[1]{Figure~\ref{#1}}
\newcommand{\chapsref}[2]{Chapters~\ref{#1}--\ref{#2}}
\newcommand{\lib}[1]{\textbf{\color{UTPBlue}\texttt{#1}}}
\newcommand{\cxg}[1]{{\color{UTPGreen}\bfseries $#1$}}
\newcommand{\cxo}[1]{{\color{UTPOrange}\bfseries $#1$}}
\newcommand{\cxr}[1]{{\color{UTPRed}\bfseries $#1$}}
\newcommand{\given}{\mathrel{|}}
\DeclareMathOperator*{\argmin}{arg\,min}
\DeclareMathOperator*{\argmax}{arg\,max}
\newcommand{\coretag}{\par\smallskip\noindent\colorbox{UTPPaleBlue}{\strut\small\bfseries\color{UTPBlue}CORE NOW}\par\smallskip}
\newcommand{\latertag}{\par\smallskip\noindent\colorbox{UTPPaleGray}{\strut\small\bfseries\color{UTPMuted}RETURN LATER}\par\smallskip}
\newcommand{\kagglesource}[2]{%
  \par\smallskip\noindent\textbf{Kaggle source:} \href{#1}{#2}\par
}

% Names and menu paths in application walkthroughs.
\newcommand{\utpapp}[1]{\textbf{\color{UTPBlue}#1}}
\newcommand{\pathmenu}[1]{\textbf{\color{UTPTeal}#1}}
\newcommand{\ptrnull}{{\scriptsize\color{UTPMuted}NULL}}

% Opt-in space beneath headings; repeated calls update the requested space.
\newcommand{\utpkeepheadings}[2]{%
  \def\utpsectionspace{#1}%
  \def\utpsubsectionspace{#2}%
  \ifdefined\utpplainsection\else
    \let\utpplainsection\section
    \let\utpplainsubsection\subsection
    \renewcommand{\section}{\keepspace{\utpsectionspace}\utpplainsection}%
    \renewcommand{\subsection}{\keepspace{\utpsubsectionspace}\utpplainsubsection}%
  \fi
}
